% =====================================================================
\section{Dispositivos de banda prohibida ancha (WBG): SiC}
\label{sec:wbg}
% =====================================================================

Todo lo anterior es silicio. Los materiales de \textbf{banda prohibida ancha} (\emph{wide band gap},
WBG) --- carburo de silicio (SiC) y nitruro de galio (GaN) --- tienen $\EG\approx\SI{3.3}{eV}$, tres
veces el del Si \diapo{25}. De ese único número salen sus dos grandes ventajas: \textbf{trabajan a
temperaturas mucho mayores} y \textbf{soportan campos diez veces mayores}.

\begin{table}[H]
\centering\small
\begin{tabular}{lcc}
\toprule
Material & $\EG$ (\si{\eV}) & $n_i$ a \SI{300}{K} (\si{cm^{-3}}) \\
\midrule
Ge & 0.66 & \num{2.4e13} \\
Si & 1.1 & \num{1.5e10} \\
GaAs & 1.4 & \num{1.8e6} \\
3C-SiC & 2.2 & 6.9 \\
6H-SiC & 3.0 & \num{2.3e-6} \\
4H-SiC & 3.26 & \num{8.2e-9} \\
GaN & 3.39 & \num{1.9e-10} \\
\bottomrule
\end{tabular}
\caption{Concentración intrínseca \diapo{25}. 3C, 6H y 4H son politipos del SiC (distintas
secuencias de apilamiento del cristal); en potencia se usa casi exclusivamente 4H-SiC por su mayor
$\EG$ y su mayor movilidad. Para el Si la diapo usa $n_i=\num{1.5e10}$; la tabla de constantes de
la materia usa \niSi, que es el valor más moderno.}
\label{tab:ni-wbg}
\end{table}

\subsection{Temperatura máxima de operación}

Una juntura funciona mientras el semiconductor sea extrínseco: $n_i\ll N_D$. Como
$n_i=\sqrt{N_cN_v}\,e^{-\EG/2kT}$ (Pierret §2.5), con $\EG$ triple el exponente es
triple.

\begin{ejemplo}[El mismo diodo a \SI{25}{\celsius} y a \SI{200}{\celsius} \diapo{26}]
Diodo de Si con $N_A=N_D=\SI{1e15}{cm^{-3}}$. A \SI{25}{\celsius}, $n_i\approx\SI{1e10}{cm^{-3}}$ y
los minoritarios son $n_i^2/N=\SI{1e5}{cm^{-3}}$. A \SI{200}{\celsius} (\SI{473}{K}), con
$N_c,N_v\propto T^{3/2}$ y $\EG(T)$ de Varshni ($\EG\approx\SI{1.08}{eV}$):
\[
  n_i(\SI{473}{K})=\sqrt{\num{3.22e19}\cdot\num{1.83e19}}\left(\tfrac{473}{300}\right)^{3/2}
  e^{-1.08/(2\cdot\num{8.61e-5}\cdot473)}\approx\SI{9e13}{cm^{-3}}\approx\SI{1e14}{cm^{-3}},
\]
y los minoritarios pasan a $\sim\SI{1e13}{cm^{-3}}$: \textbf{ocho órdenes de magnitud} más. La
corriente de saturación ($\propto n_i^2$) crece en la misma proporción y, poco más arriba, $n_i$
alcanza al dopaje y la juntura deja de existir.
\end{ejemplo}

\begin{figure}[H]
\centering
\begin{tikzpicture}
\begin{axis}[
  width=0.82\textwidth, height=6.6cm,
  ymode=log, xmin=0, xmax=600, ymin=1e-10, ymax=1e18,
  xlabel={$T$ (\si{\celsius})}, ylabel={$n_i$ (\si{cm^{-3}})},
  ytick={1e-10,1e-6,1e-2,1e2,1e6,1e10,1e14,1e18},
  grid=major, grid style={gris!20},
  legend pos=south east, legend style={font=\small},
  domain=0:600, samples=80,
  every axis plot/.append style={line width=1.1pt},
]
  \addplot[rojo] {sqrt(3.22e19*1.83e19)*((x+273)/300)^1.5*exp(-(1.17-4.73e-4*(x+273)^2/(x+273+636))/(2*8.61e-5*(x+273)))};
  \addlegendentry{Si}
  \addplot[azul] {sqrt(1.7e19*2.5e19)*((x+273)/300)^1.5*exp(-3.26/(2*8.61e-5*(x+273)))};
  \addlegendentry{4H-SiC}
  \addplot[verde!60!black] {sqrt(2.3e18*4.6e19)*((x+273)/300)^1.5*exp(-3.39/(2*8.61e-5*(x+273)))};
  \addlegendentry{GaN}
  \addplot[gris,dashed,line width=0.8pt] {1e15};
  \node[font=\small,anchor=north west] at (axis cs:300,1e15) {$N_D=\SI{1e15}{cm^{-3}}$ (arrastre)};
  \draw[rojo,dotted,line width=0.9pt] (axis cs:280,1e-10) -- (axis cs:280,1e15);
  \node[rojo,font=\small,anchor=east] at (axis cs:278,1e-4) {Si: $\approx\SI{280}{\celsius}$};
\end{axis}
\end{tikzpicture}
\caption{$n_i(T)$ calculada con $\sqrt{N_cN_v}(T/300)^{3/2}e^{-\EG/2kT}$ \diapo{27}. Para Si se usan
$N_c$, $N_v$ de la tabla de constantes y $\EG(T)$ de Varshni; para 4H-SiC y GaN, valores de
literatura con $\EG$ constante. El Si con $N_D=\SI{1e15}{cm^{-3}}$ se vuelve intrínseco a
$\sim\SI{280}{\celsius}$; el SiC y el GaN recién por encima de \SI{1000}{\celsius} (en la práctica los
limita el encapsulado, a $\sim\SIrange{200}{250}{\celsius}$).}
\label{fig:ni-T}
\end{figure}

\subsection{Comparación de materiales}

\begin{table}[H]
\centering\small
\begin{tabular}{lccccccc}
\toprule
Material & $\EG$ & $\mu_n$ & $\vsat$ & $\Ecr$ & $\lambda$ & $\epsilon_r$ & BFOM \\
 & (\si{\eV}) & (\si{cm^2/V.s}) & (\si{cm/s}) & (\si{V/cm}) & (\si{W/cm.K}) & & (rel.\ Si) \\
\midrule
Si      & 1.12 & 1350 & \num{1e7}   & \num{3e5}   & 1.5 & 11.8 & 1 \\
GaAs    & 1.42 & 8500 & \num{2e7}   & \num{4e5}   & 0.5 & 13.1 & 17 \\
3C-SiC  & 2.2  & 900  & \num{2e7}   & \num{1.2e6} & 4.5 & 9.7  & 35 \\
6H-SiC  & 3.0  & 370  & \num{2e7}   & \num{2.4e6} & 4.5 & 9.7  & 115 \\
4H-SiC  & 3.26 & 720  & \num{2e7}   & \num{2e6}   & 4.5 & 9.7  & 130 \\
GaN     & 3.39 & 900  & \num{2.5e7} & \num{3.3e6} & 1.3 & 9.0  & 680 \\
\bottomrule
\end{tabular}
\caption{Propiedades de los materiales \diapo{29} (se agregó $\epsilon_r$ y la figura de mérito de
Baliga $\epsilon\mu_n\Ecr^3$ relativa al Si, calculada con esos valores). En el GaN el valor de
$\mu_n$ es de volumen; en el canal 2D de un HEMT llega a \SIrange{1500}{2000}{cm^2/V.s}
(sección~\ref{sec:gan}).}
\label{tab:materiales}
\end{table}

\begin{figure}[H]
\centering
\begin{tikzpicture}[scale=1]
  % ejes: (nombre, max, ángulo)
  \def\R{2.6}
  \foreach \a/\lab in {90/{Campo crítico\\(MV/cm, máx.\ 4)},18/{Conductividad\\térmica (W/cm\,K, máx.\ 5)},
                     -54/{T. de fusión\\(\num{1000}\,\si{\celsius}, máx.\ 3)},-126/{Velocidad de sat.\\($10^7$\,cm/s, máx.\ 3)},
                     162/{Energía del gap\\(eV, máx.\ 4)}} {
    \draw[gris!50] (0,0) -- (\a:\R);
    \node[font=\scriptsize,align=center] at (\a:{\R+0.85}) {\lab};
  }
  \foreach \f in {0.25,0.5,0.75,1} \draw[gris!30] (90:\f*\R) -- (18:\f*\R) -- (-54:\f*\R) -- (-126:\f*\R) -- (162:\f*\R) -- cycle;
  % Si: Ec .3/4, k 1.5/5, Tf 1.42/3, vs 1/3, Eg 1.12/4
  \draw[naranja,line width=1.2pt,fill=naranja,fill opacity=0.12]
    (90:{0.3/4*\R}) -- (18:{1.5/5*\R}) -- (-54:{1.42/3*\R}) -- (-126:{1/3*\R}) -- (162:{1.12/4*\R}) -- cycle;
  % 4H-SiC: Ec 2, k 4.5, Tf 2.8, vs 2, Eg 3.26
  \draw[azul,line width=1.2pt,fill=azul,fill opacity=0.12]
    (90:{2/4*\R}) -- (18:{4.5/5*\R}) -- (-54:{2.8/3*\R}) -- (-126:{2/3*\R}) -- (162:{3.26/4*\R}) -- cycle;
  % GaN: Ec 3.3, k 1.3, Tf 2.5, vs 2.5, Eg 3.39
  \draw[rojo,line width=1.2pt,fill=rojo,fill opacity=0.12]
    (90:{3.3/4*\R}) -- (18:{1.3/5*\R}) -- (-54:{2.5/3*\R}) -- (-126:{2.5/3*\R}) -- (162:{3.39/4*\R}) -- cycle;
  % leyenda
  \begin{scope}[xshift=6.6cm,yshift=-0.6cm]
    \draw[naranja,line width=1.2pt] (0,0) -- (0.5,0) node[right,black,font=\small] {Si};
    \draw[azul,line width=1.2pt] (0,-0.45) -- (0.5,-0.45) node[right,black,font=\small] {4H-SiC};
    \draw[rojo,line width=1.2pt] (0,-0.9) -- (0.5,-0.9) node[right,black,font=\small] {GaN};
  \end{scope}
\end{tikzpicture}
\caption{Comparación de materiales \diapo{28}. Campo crítico y gap $\Rightarrow$ \textbf{alta
tensión}; conductividad térmica y temperatura de fusión $\Rightarrow$ \textbf{alta temperatura}
(el SiC disipa el calor tres veces mejor que el Si); velocidad de saturación $\Rightarrow$ \textbf{alta
frecuencia} (el GaN es el mejor).}
\label{fig:radar}
\end{figure}

\subsection{Dos uniones que soportan la misma tensión}

\diapo{30} Se construyen dos diodos $p^+n^-$, uno de Si y otro de SiC, para la misma $\VB$. Con las
fórmulas de la sección~\ref{sec:deriva}:
\[
  \frac{N_{D,\mathrm{SiC}}}{N_{D,\mathrm{Si}}}=\frac{\epsilon_{\mathrm{SiC}}\mathcal{E}_{c,\mathrm{SiC}}^2}{\epsilon_{\mathrm{Si}}\mathcal{E}_{c,\mathrm{Si}}^2},\qquad
  \frac{W_{\mathrm{Si}}}{W_{\mathrm{SiC}}}=\frac{\mathcal{E}_{c,\mathrm{SiC}}}{\mathcal{E}_{c,\mathrm{Si}}},\qquad
  \frac{R_{\mathrm{Si}}}{R_{\mathrm{SiC}}}=\frac{(\epsilon\mu_n\Ecr^3)_{\mathrm{SiC}}}{(\epsilon\mu_n\Ecr^3)_{\mathrm{Si}}}.
\]
Y si en cambio se construyen con la \emph{misma resistencia}, como $\Ronsp=4\VB^2/(\epsilon\mu_n\Ecr^3)$:
\[
  \frac{V_{B,\mathrm{SiC}}}{V_{B,\mathrm{Si}}}=\sqrt{\frac{(\epsilon\mu_n\Ecr^3)_{\mathrm{SiC}}}{(\epsilon\mu_n\Ecr^3)_{\mathrm{Si}}}}
  \qquad\text{\diapo{30}}.
\]

\begin{ejemplo}[Si vs.\ 4H-SiC para \SI{600}{V}]
Con los valores de la tabla~\ref{tab:materiales}:
\begin{center}\small
\begin{tabular}{lccc}
\toprule
 & $N_D$ & $W$ & $\Ronsp$ \\
\midrule
Si     & \SI{4.9e14}{cm^{-3}} & \SI{40}{\micro m} & \SI{38}{m\ohm.cm^2} \\
4H-SiC & \SI{1.8e16}{cm^{-3}} & \SI{6}{\micro m}  & \SI{0.29}{m\ohm.cm^2} \\
\midrule
cociente & 36.5 & 6.7 & 130 \\
\bottomrule
\end{tabular}
\end{center}
La región de arrastre del SiC se dopa 36 veces más y es casi 7 veces más delgada; su resistencia es
130 veces menor. A igual resistencia, el SiC bloquea $\sqrt{130}\approx11$ veces más tensión.
\end{ejemplo}

\begin{ojo}[Los números de la diapo 30]
La diapo da $N_{D,\mathrm{SiC}}=26.14\,N_{D,\mathrm{Si}}$, $R_{N,\mathrm{SiC}}=R_{N,\mathrm{Si}}/92.94$ y
un factor $9.6$ en tensión ($=\sqrt{92.94}$, consistente con la $R$). Con los valores de su propia
tabla (diapo~29: $\Ecr=\SI{2e6}{V/cm}$ y $\mu_n=720$ para 4H-SiC) se obtienen 36.5, 130 y 11.4. La
diferencia viene de los parámetros elegidos (el $\Ecr$ del SiC se cita entre 2 y \SI{3}{MV/cm}, y
depende del dopaje); lo robusto es el orden de magnitud: \textbf{$\sim10$ veces más tensión, o
$\sim100$ veces menos resistencia}.
\end{ojo}

\subsection{Diodos de SiC}

\diapo{31} Tres tipos:
\begin{itemize}
  \item \textbf{Schottky} (hasta $\approx\SIrange{600}{1200}{V}$, hoy \SI{1700}{V}): unipolar, sin
  carga almacenada ni recuperación inversa. Es el producto de SiC más maduro: reemplaza al diodo
  ultrarrápido de Si en fuentes con corrección de factor de potencia e inversores.
  \item \textbf{PN y PiN} (varios kV): bipolares, modulan la región de arrastre, pero con una caída
  de codo de casi \SI{3}{V} (ver ejemplo).
  \item \textbf{Mixtos} (JBS/MPS, tensiones intermedias): un contacto Schottky con islas $p^+$
  intercaladas. En inversa, las zonas desiertas de las islas se juntan y \emph{apantallan} la
  interfaz Schottky (menos fuga, más tensión); en directa conduce el Schottky y, ante sobrecorrientes,
  las islas inyectan como un PiN.
\end{itemize}
Regla general: \textbf{cuanto más Schottky, más rápido y menor caída en conducción, pero menor tensión
inversa y más fugas; cuanto más bipolar (PN), lo contrario} \diapo{31}.

\begin{ejemplo}[Por qué un diodo PN de SiC cae \SI{3}{V}]
La corriente de un diodo es $J=J_0(e^{V/V_T}-1)$ con $J_0\propto n_i^2$ (hoja de fórmulas,
$I_S$ \enhojatexto). Para la misma densidad de corriente,
\[
  V_{\mathrm{SiC}}-V_{\mathrm{Si}}=V_T\ln\frac{n_{i,\mathrm{Si}}^2}{n_{i,\mathrm{SiC}}^2}
  =2\cdot\SI{0.026}{V}\cdot\ln\frac{\num{1e10}}{\num{8.2e-9}}\approx\SI{2.2}{V}.
\]
Si un diodo de Si conduce a \SI{0.7}{V}, uno de 4H-SiC lo hace a $\approx\SI{2.9}{V}$: el gap grande
que permite bloquear tanto también encarece cada ampere en conducción bipolar. Por eso en SiC se
prefieren los dispositivos unipolares (Schottky, MOSFET, JFET) hasta varios kV.
\end{ejemplo}

\subsection{JFET de SiC}

\diapo{32} El JFET fue el primer transistor de SiC comercial, porque esquiva el problema del óxido
(abajo).
\begin{center}\small
\begin{tabularx}{0.95\textwidth}{XX}
\toprule
\textbf{Ventajas} & \textbf{Desventajas} \\
\midrule
Conducción unipolar (sin degradación bipolar). & Es \emph{normalmente cerrado} (\emph{normally-on}):
conduce con $V_{GS}=0$. \\
Rápido. & Se apaga con tensión de compuerta negativa. \\
Se fabrica sobre obleas $n$; estructura simple. & \\
\bottomrule
\end{tabularx}
\end{center}

\begin{ojo}[``Abierto'' y ``cerrado'']
En las diapositivas \emph{normalmente cerrado} significa lo que en inglés se dice
\emph{normally-on}: la llave está cerrada, es decir \textbf{conduce} sin tensión de compuerta (JFET,
HEMT de deplexión). \emph{Normalmente abierto} = \emph{normally-off} (MOSFET de enriquecimiento). Es la
convención de llaves y relés; no confundir ``abierto'' con ``canal abierto''.
\end{ojo}

\paragraph{¿Abiertos o cerrados? \diapo{33}} En el convertidor de la diapo~33 (puente de JFETs), al
energizar el circuito, antes de que el control genere la tensión negativa de compuerta, \textbf{los
JFET están todos conduciendo}: las dos llaves de una rama cortocircuitan la fuente (\emph{shoot-through}).
Un dispositivo normally-on exige una secuencia de arranque cuidadosa o una solución de circuito:

\begin{figure}[H]
\centering
\begin{circuitikz}
  \draw (0,3.0) node[njfet] (J) {};
  \draw (0,0.8) node[nmos] (M) {};
  \draw (J.D) -- ++(0,0.4) node[above] {D};
  \draw (J.S) -- (M.D);
  \draw (M.S) -- ++(0,-0.3) coordinate (s) -- ++(0,-0.3) node[below] {S};
  \draw (J.G) -- ++(-0.5,0) -- ++(0,-2.4) coordinate (gj) |- (s);
  \draw (M.G) -- ++(-0.35,0) node[left] {G};
  \node[right,font=\small] at (J.east) {JFET SiC (normally-on, alta tensión)};
  \node[right,font=\small] at (M.east) {MOSFET Si (normally-off, baja tensión)};
\end{circuitikz}
\caption{Cascodo JFET-SiC $+$ MOSFET-Si \diapo{32}. La compuerta del JFET está atada a la fuente del
conjunto. Con el MOSFET encendido, $V_{GS,\mathrm{JFET}}\approx0$ y el JFET conduce. Con el MOSFET
apagado, su $V_{DS}$ sube hasta que $V_{GS,\mathrm{JFET}}=-V_{DS,\mathrm{MOS}}$ alcanza la tensión de
estrangulamiento del JFET ($\sim\SI{-10}{V}$), que se apaga y toma el resto de la tensión. El
conjunto se comporta como una llave \emph{normally-off} controlada por la compuerta del MOSFET, que
sólo tiene que bloquear unas decenas de volts.}
\label{fig:cascodo}
\end{figure}

\subsection{MOSFET de SiC}

El SiC tiene una ventaja única entre los WBG: como el Si, \textbf{forma $\mathrm{SiO_2}$ por oxidación
térmica}, así que admite MOSFET de compuerta aislada. Pero \diapo{34}:
\begin{itemize}
  \item \textbf{Baja movilidad en el canal}: la interfaz $\mathrm{SiC/SiO_2}$ tiene muchas más trampas y
  cargas que la de Si (en parte por el carbono residual). Los electrones del canal tienen
  \SIrange{20}{50}{cm^2/V.s} frente a 720--\SI{900}{cm^2/V.s} en volumen: en un MOSFET de SiC la
  resistencia del canal pesa tanto como la de la región de arrastre.
  \item \textbf{El óxido es clave} y es difícil de lograr con buena calidad y confiabilidad (deriva de
  $V_{TH}$ con la temperatura y el uso).
  \item \textbf{No reemplazan directamente} a los MOSFET de Si: se excitan con $V_{GS}\approx+15$ a
  $+20$\,V para encender (vs.\ \SIrange{10}{12}{V}) y con $-3$ a $-5$\,V para asegurar el apagado.
\end{itemize}
La figura de la diapo~34 es un MOSFET \emph{lateral} (didáctico); los MOSFET de SiC comerciales son
verticales (DMOS o \emph{trench}) como el de la figura~\ref{fig:vdmos}.

\begin{ejemplo}[MOSFET SiC vs.\ MOSFET Si de \SI{900}{V} \diapo{35}]
\begin{center}\small
\begin{tabular}{lccc}
\toprule
 & $\RDSon$ & $\Qg$ & $\RDSon\Qg$ \\
\midrule
SiC (CREE C3M0065090J) & \SI{65}{m\ohm} & \SI{30}{nC} & \SI{1.95}{\ohm.nC} \\
Si superjuntura (Infineon IPP90R340C3) & \SI{340}{m\ohm} & \SI{94}{nC} & \SI{32.0}{\ohm.nC} \\
\bottomrule
\end{tabular}
\end{center}
Cociente: $32.0/1.95=16.4$, ``aproximadamente 16 veces mejor el SiC'' \diapo{35}. También las
capacidades: $C_{iss}=\SI{660}{pF}$ vs.\ \SI{2400}{pF}, $C_{oss}=\SI{60}{pF}$ vs.\ \SI{120}{pF}. El SiC
tiene cinco veces menos resistencia \emph{y} tres veces menos carga de compuerta a la vez.
\end{ejemplo}
